\begin{mistake}{2026-04-13}{07}

\begin{problemcontent}
设 \(A\) 为 3 阶矩阵，已知 \(\boldsymbol{\alpha}_1, \boldsymbol{\alpha}_2, \boldsymbol{\alpha}_3\) 是线性无关的 3 维列向量，且满足：
\(A\boldsymbol{\alpha}_1 = \boldsymbol{\alpha}_1 + \boldsymbol{\alpha}_2 + \boldsymbol{\alpha}_3\)，
\(A\boldsymbol{\alpha}_2 = 2\boldsymbol{\alpha}_2 + \boldsymbol{\alpha}_3\)，
\(A\boldsymbol{\alpha}_3 = 2\boldsymbol{\alpha}_2 + 3\boldsymbol{\alpha}_3\)。
(1) 求矩阵 \(B\)，使得 \(A[\boldsymbol{\alpha}_1, \boldsymbol{\alpha}_2, \boldsymbol{\alpha}_3] = [\boldsymbol{\alpha}_1, \boldsymbol{\alpha}_2, \boldsymbol{\alpha}_3]B\)；
(2) 求矩阵 \(A\) 的特征值；
(3) 求可逆矩阵 \(P\)，使得 \(P^{-1}AP\) 为对角矩阵。
\end{problemcontent}

\begin{solutioncontent}
（1）令 \(C = [\boldsymbol{\alpha}_1, \boldsymbol{\alpha}_2, \boldsymbol{\alpha}_3]\)。由已知等式和矩阵乘法规则得：
\[
AC = [A\boldsymbol{\alpha}_1, A\boldsymbol{\alpha}_2, A\boldsymbol{\alpha}_3] = [\boldsymbol{\alpha}_1, \boldsymbol{\alpha}_2, \boldsymbol{\alpha}_3]
\begin{pmatrix}
1 & 0 & 0 \\
1 & 2 & 2 \\
1 & 1 & 3
\end{pmatrix}
\]
由于 \(\boldsymbol{\alpha}_1, \boldsymbol{\alpha}_2, \boldsymbol{\alpha}_3\) 线性无关，矩阵 \(C\) 可逆，故
\(B = \begin{pmatrix} 1 & 0 & 0 \\ 1 & 2 & 2 \\ 1 & 1 & 3 \end{pmatrix}\)。

（2）由 \(AC = CB\) 得 \(C^{-1}AC = B\)，因此 \(A \sim B\)，两者特征值相同。
\[
\begin{aligned}
|\lambda E - B| &= \begin{vmatrix}
\lambda - 1 & 0 & 0 \\
-1 & \lambda - 2 & -2 \\
-1 & -1 & \lambda - 3
\end{vmatrix} \\
&= (\lambda - 1)[(\lambda - 2)(\lambda - 3) - 2] \\
&= (\lambda - 1)^2(\lambda - 4) = 0
\end{aligned}
\]
解得 \(B\) 的特征值，即 \(A\) 的特征值为 \(\lambda_1 = \lambda_2 = 1, \lambda_3 = 4\)。

（3）先求 \(B\) 的特征向量。
当 \(\lambda_{1,2} = 1\) 时，解 \((E - B)\boldsymbol{x} = \boldsymbol{0}\)，得线性无关的特征向量 \(\boldsymbol{\xi}_1 = (-1, 1, 0)^T\), \(\boldsymbol{\xi}_2 = (-2, 0, 1)^T\)。
当 \(\lambda_3 = 4\) 时，解 \((4E - B)\boldsymbol{x} = \boldsymbol{0}\)，得特征向量 \(\boldsymbol{\xi}_3 = (0, 1, 1)^T\)。
令 \(Q = [\boldsymbol{\xi}_1, \boldsymbol{\xi}_2, \boldsymbol{\xi}_3] = \begin{pmatrix} -1 & -2 & 0 \\ 1 & 0 & 1 \\ 0 & 1 & 1 \end{pmatrix}\)，则 \(Q^{-1}BQ = \text{diag}(1, 1, 4)\)。
将 \(B = C^{-1}AC\) 代入得 \(Q^{-1}(C^{-1}AC)Q = (CQ)^{-1}A(CQ) = \text{diag}(1, 1, 4)\)。
故所求可逆矩阵 \(P = CQ\)，展开得：
\[
P = [-\boldsymbol{\alpha}_1 + \boldsymbol{\alpha}_2, \quad -2\boldsymbol{\alpha}_1 + \boldsymbol{\alpha}_3, \quad \boldsymbol{\alpha}_2 + \boldsymbol{\alpha}_3]
\]
\end{solutioncontent}

\mistakeknowledge{
\begin{itemize}
 \item \textbf{核心公式/定理}：分块矩阵乘法法则；相似矩阵定义 \(C^{-1}AC = B\)，相似矩阵具相同的特征值。
 \item \textbf{易错点}：求坐标矩阵 \(B\) 时，易错把方程系数按行写入（应按列写）；求 \(P\) 时易直接把对角化 \(B\) 的变换矩阵 \(Q\) 当作答案，遗漏左乘矩阵 \(C\)。
 \item \textbf{关联知识}：若 \(B\boldsymbol{x}=\lambda\boldsymbol{x}\) 且 \(C^{-1}AC=B\)，则 \(C^{-1}AC\boldsymbol{x}=\lambda\boldsymbol{x}\)，进而 \(A(C\boldsymbol{x})=\lambda(C\boldsymbol{x})\)，这证明了 \(C\boldsymbol{x}\) 是 \(A\) 的特征向量。
\end{itemize}}

\end{mistake}
